\chapter{Project Plan}
\section{Overview}
This chapter provides a plan for the remaining work to be completed for the 
project. Firstly, we summarise our finding so far, and comment on what was and
wasn't achieved. Then, we outline the methodology for the remaining work and how 
this plan differs from the initial proposal. The next section will 
describe the 
timeline for the remaining work. Then, the plan for the management of 
the project stakeholders and risks will be outlined. Finally a reflection of 
how the project has contributed to the learning objectives of the AMME5520 and 
ELEC3305 units will be provided.

\section{Summary of Findings}
\todo{Summarise the finds at a high level. State the best/recommended configuration.}
\todo{Link together all the sections of the discussion into main takaways. Kinda like an executive summary.}
A modular radar pipeline for small-UAV detect-and-avoid was successfully designed 
and evaluated.

CFAR was validated, with achieved false alarm rates closely matching desired 
values.

Implementation choice strongly affected runtime, latency, and throughput across 
hardware platforms.

Hardware context changed which implementation was most suitable, supporting 
hardware-aware pipeline design.

\section{Industry Partner Outcomes}
\todo{Without sharing anything sensitive, detail how the research done benefits mission systems in their project goals.}
\todo{Document tests and downstream work that my work enabled.}
\todo{Discuss how this technology will after the industry outside of Mission Systems too.}
Delivered practical benchmarking evidence to guide embedded radar software 
design choices.

Helped identify implementation trade-offs relevant to Mission Systems’ UAV radar 
pipeline.

Produced a reusable experimental and profiling framework for future internal 
development.

Contributed design guidance with relevance beyond the immediate industry project.

\section{Future Work}
\todo{What about our research was successful, what wasn't successful}
\todo{Discuss the limitations of the set of computing hardware that we were able to explore}
\todo{We only explored one radar device}
\todo{Talk about algorithms that were not explored}
\todo{Compare to other sensing systems}
\todo{Sensor Fusion}

(See Addressing Knowledge Gaps in Chapter 6.)